% Part: history
% Chapter: biographies
% Section: alan-turing

\documentclass[../../../include/open-logic-section]{subfiles}

\begin{document}

\olfileid[id]{his}{bio}{tur}

\olsection{Alan Turing}

Alan Turing lahir di Maida Vale, London, pada 23 Juni 1912. Ia dianggap
sebagai bapak ilmu komputer teoretis. Minat Turing terhadap ilmu-ilmu fisik
dan matematika dimulai sejak usia muda. Namun, ketika ia masih anak-anak,
minatnya tidak terwakili dengan baik di sekolah-sekolahnya, yang lebih
menekankan sastra dan karya-karya klasik. Akibatnya, prestasinya di sekolah
buruk dan banyak gurunya menegurnya.

\olphoto{turing-alan}{Alan Turing}

Turing menempuh pendidikan sarjana di King's College, Cambridge, dan
mempelajari matematika. Pada 1936, Turing mengembangkan apa yang sekarang
disebut mesin Turing sebagai upaya mendefinisikan secara tepat pengertian
fungsi yang dapat dihitung dan membuktikan ketakterputusan masalah
keputusan. Alonzo Church mendahuluinya memperoleh hasil tersebut melalui
kalkulus lambda ciptaannya sendiri. Makalah Turing tetap diterbitkan dengan
merujuk hasil Church. Church mengundang Turing ke Princeton, tempat ia
tinggal pada 1936--1938 dan memperoleh gelar doktor di bawah bimbingan
Church.

Meskipun berminat pada logika, minat awal Turing terhadap ilmu-ilmu fisik
tetap kuat. Keterampilan praktisnya digunakan selama ia bertugas di
departemen kriptanalisis Britania di Bletchley Park pada Perang Dunia~II\@.
Turing menjadi tokoh sentral dalam pemecahan sandi komunikasi Angkatan Laut
Jerman---kode Enigma. Keahlian Turing dalam statistika dan kriptografi,
bersama penggunaan mesin elektronik, memungkinkan tim tersebut memecahkan
kode dengan menciptakan mesin pengurai sandi yang disebut ``bombe''.
Gagasannya juga membantu penciptaan Colossus, komputer elektronik pertama
di dunia yang dapat diprogram, yang juga digunakan di Bletchley Park untuk memecahkan
sandi Lorenz Jerman.

Turing adalah seorang gay. Meskipun demikian, pada 1942 ia melamar Joan Clarke,
salah seorang rekan satu timnya di Bletchley Park, tetapi kemudian
membatalkan pertunangan itu dan mengaku kepadanya bahwa ia homoseksual. Ia
mempunyai beberapa kekasih sepanjang hidupnya, meskipun tindakan homoseksual
pada waktu itu merupakan tindak pidana di Britania Raya. Pada 1952, rumah
Turing dibobol oleh seorang teman kekasihnya saat itu. Ketika melaporkannya
kepada polisi, Turing mengakui bahwa ia menjalani hubungan homoseksual
karena mengira pemerintah sedang menuju legalisasi tindakan homoseksual.
Anggapan itu tidak benar, dan ia didakwa melakukan perbuatan tidak senonoh
berat. Alih-alih dipenjara, Turing memilih terapi hormon yang menurunkan
libido. Turing ditemukan meninggal pada 8 Juni 1954 akibat overdosis
sianida---kemungkinan besar bunuh diri. Ratu Elizabeth~II memberinya
pengampunan kerajaan pada 2013.

\begin{reading}
Untuk biografi Alan Turing yang menyeluruh, lihat \citet{Hodges2014}.
Kehidupan dan karya Turing mengilhami drama \emph{Breaking the Code}, yang
diproduksi untuk televisi pada 1996 dengan Derek Jacobi sebagai Turing.
\emph{The Imitation Game}, film nominasi Academy Award yang dibintangi
Benedict Cumberbatch dan Keira Knightley, juga didasarkan secara longgar pada
kehidupan Alan Turing dan masanya di Bletchley Park
\citep{Imitation2014}.

\citet{Radiolab2012} mempunyai beberapa siniar tentang kehidupan dan karya
Turing. Dokumenter BBC Horizon \emph{The Strange Life and Death of
  Dr.~Turing} tersedia untuk ditonton daring \citep{Sykes1992}.
\citep{Theelen2012} merupakan video singkat sebuah Mesin Turing LEGO yang
berfungsi---dibuat untuk menghormati seratus tahun kelahiran Turing pada
2012.

Makalah asli Turing mengenai mesin Turing dan masalah keputusan ialah
\citet{Turing1937}.
\end{reading}

\end{document}
